% Section 4.2, from lectures/L04/README.md: the goals, the evaluation questions, and the next
% lecture.

\section{Review}
\label{c4:sec:review}

\needspace{6\baselineskip}
\subsection*{What you should be able to do now}
After this chapter, you should:
\begin{itemize}
  \item Explain what an ADC does and why its result is a discrete number, not a continuous value.
  \item Be able to compute the voltage a given ADC reading corresponds to, and vice versa.
  \item Be able to configure and read from a chosen ADC channel.
  \item Understand why the ADC needs its own clock prescaler, separate from the CPU clock.
\end{itemize}

\needspace{6\baselineskip}
\subsection*{Questions to test yourself}
You should be able to answer:
\begin{enumerate}
  \item At a 5\,V reference, what analog voltage does an ADC reading of \code{0}, \code{511}, and
    \code{1023} each correspond to?
  \item Why does the code wait on the \code{ADIF} flag instead of reading \code{ADC} immediately
    after starting a conversion?
  \item What is ``duty cycle,'' and how does the software-PWM example use it to control perceived
    LED brightness?
\end{enumerate}

\needspace{6\baselineskip}
\subsection*{Looking ahead}
Chapter~\ref{c5:ch} is about timers:
\begin{itemize}
  \item Timers: replacing blocking \code{_delay_ms()} calls with non-blocking, interrupt-driven
    timing.
  \item Normal Mode vs.\ CTC Mode, and the prescaler math behind both.
\end{itemize}
